%% ============================================================
%%  Chapter 3 — Injection (example category)
%%  Findings: F-001 to F-00N
%%
%%  Rename this file/chapter to whatever vulnerability category
%%  applies (e.g. "Authentication", "Authorization", "Secrets \&
%%  Credentials"), and copy the F-001 block below once per finding.
%% ============================================================
\chapter{Injection}
\label{ch:example_category_one}

\begin{warningbox}
  \textbf{Systemic context:} \fillme{If this category has a systemic root cause -- e.g. a shared
  data-access helper that never parameterises input -- describe it here once, so each finding
  below can just reference it instead of repeating the explanation.}
\end{warningbox}

% ---- F-001 ------------------------------------------------------------------
\section{F-001 --- \fillme{Short, specific finding title}}
\label{sec:F-001}

\begin{findingcritical}{F-001}{\fillme{Short, specific finding title}}
  \findmeta{Severity}{\Critical} \findmeta{Confidence}{\Confirmed}

  \bigskip
  \findsection{Location}
  \cls{\fillme{package/Class.java}}; reachable from \fillme{endpoint(s) or entry point(s)}.

  \findsection{Description}
  \fillme{Explain, in plain language, what untrusted input reaches the vulnerable sink and why
  it is not safely handled.}

  \begin{lstlisting}[style=sql, caption={\fillme{Illustrative injectable query -- no real schema/data}}]
-- minimal illustrative SQL showing the pattern -- replace with your own, not real table/column names
  \end{lstlisting}

  \findsection{Proof of Concept (illustrative --- not executed)}
  \begin{lstlisting}[style=http, caption={\fillme{PoC title}}]
# illustrative request only -- replace with your own, no real hostnames/tokens/credentials
  \end{lstlisting}

  \findsection{Impact}
  \fillme{One or two sentences: what an attacker can concretely do as a result.}

  \findsection{Remediation}
  \fillme{The specific fix, e.g. "use parameterised queries / bind variables everywhere this
  pattern occurs."}
\end{findingcritical}

% ---- F-002 (medium-severity example) ----------------------------------------
\section{F-002 --- \fillme{Second finding title in this category}}
\label{sec:F-002}

\begin{findingmedium}{F-002}{\fillme{Second finding title}}
  \findmeta{Severity}{\Medium} \findmeta{Confidence}{\Suspected}

  \bigskip
  \findsection{Location} \cls{\fillme{package/Class.java}}

  \findsection{Description}
  \fillme{Describe the weaker/partial version of the same issue class.}

  \findsection{Impact} \fillme{...}

  \findsection{Remediation} \fillme{...}
\end{findingmedium}
